%% =====================================================================
%%  core/lib/partes-comunes.tex  --  Partes compartidas por TODOS los
%%                                   tipos documentales
%% ---------------------------------------------------------------------
%%  Un solo lugar para lo que antes se repetia en cada estructura:
%%
%%    \bloquefirma[espacio]   firma centrada + sello de consigna
%%    \adjunto{...}           "Adj.: ..."
%%    \concopia{...}          "c.c.: ..."
%%    \selloconsigna          el sello inclinado (puno + consigna + puno)
%%
%%  El SELLO es el ARTE de la consigna (bloques rojo/amarillo tipo
%%  cartel + puno a la izquierda), generado por scripts/generar-sello.py
%%  en assets/logos/sello-consigna.png, colocado a la DERECHA del
%%  bloque de firma. Si la consigna cambia, editar el script y
%%  regenerar el PNG (la frase va horneada en el arte).
%%
%%    \setmeta{sello}{no|si}       interruptor por METADATOS (documento)
%%    \sinsello / \consello        desactivar / reactivar (solo colgante)
%%    \setsellotamano{2.2cm}       alto del arte
%%    \setselloangulo{0}           inclinacion opcional en grados
%%    \setselloposicion{dim}       distancia margen izq. -> sello
%%    \setselloaltura{dim}         elevacion sobre la linea de firma
%% =====================================================================

\RequirePackage{eso-pic}

%% --- Sello de consigna: bandera y ajustes finos ----------------------
\newif\ifdespachosello
\despachosellotrue
\newcommand{\consello}{\despachosellotrue}
\newcommand{\sinsello}{\despachosellofalse}

%% --- Sello por METADATOS (interruptor a nivel documento) -------------
%%  \setmeta{sello}{no} en el metadata.tex del documento apaga TODAS
%%  las variantes del sello (el colgante de \bloquefirma Y el de
%%  \selloalfondo); \setmeta{sello}{si} lo fuerza. Sin la clave, el
%%  comportamiento por defecto (con sello).
%%  OJO: es una bandera DISTINTA de \sinsello, porque el patron
%%  multipagina usa \sinsello + \selloalfondo (ahi \sinsello significa
%%  "no colgar de la firma", no "documento sin sello").
\newif\ifdespachosellodoc
\despachosellodoctrue
\AtBeginDocument{%
  \ifcsname despacho@meta@sello\endcsname
    \edef\despachoSelloMeta{\csname despacho@meta@sello\endcsname}%
    \def\despachoSelloMetaNo{no}%
    \def\despachoSelloMetaSi{si}%
    \ifx\despachoSelloMeta\despachoSelloMetaNo
      \despachosellodocfalse
      \sinsello
    \fi
    \ifx\despachoSelloMeta\despachoSelloMetaSi
      \despachosellodoctrue
      \consello
    \fi
  \fi
}

\newcommand{\despachoSelloAngulo}{0}
\newcommand{\setselloangulo}[1]{\renewcommand{\despachoSelloAngulo}{#1}}

%%  El sello va DEBAJO del bloque de firma (bajo el cargo), ALINEADO
%%  AL MARGEN DERECHO.
\newlength{\despachoSelloTam}
\setlength{\despachoSelloTam}{1.20cm}
\newcommand{\setsellotamano}[1]{\setlength{\despachoSelloTam}{#1}}

%%  Separacion desde el margen derecho (0 = al ras del margen)
\newlength{\despachoSelloPos}
\setlength{\despachoSelloPos}{0cm}
\newcommand{\setselloposicion}[1]{\setlength{\despachoSelloPos}{#1}}

%%  Separacion vertical entre el cargo y el sello
\newlength{\despachoSelloAlt}
\setlength{\despachoSelloAlt}{0.30cm}
\newcommand{\setselloaltura}[1]{\setlength{\despachoSelloAlt}{#1}}

%% --- El sello --------------------------------------------------------
%%  El arte completo (puno + consigna) en un solo PNG transparente; el
%%  archivo lo fija el perfil con \setselloimagen (neutro: sello.png).
\newcommand{\despachoSelloImagen}{sello.png}
\newcommand{\setselloimagen}[1]{\renewcommand{\despachoSelloImagen}{#1}}
%%  Retrato de cuerpo para la contraportada del informe-sr (\setretratocuerpo) y de busto para la
%%  presentación (\setretratobusto, DOC7 2026-09-20; por defecto, el mismo archivo que el de cuerpo).
\providecommand{\despachoRetratoBusto}{\despachoRetratoCuerpo}
\providecommand{\setretratobusto}[1]{\renewcommand{\despachoRetratoBusto}{#1}}
\newcommand{\despachoRetratoCuerpo}{retrato-cuerpo.png}
\newcommand{\setretratocuerpo}[1]{\renewcommand{\despachoRetratoCuerpo}{#1}}
\newcommand{\selloconsigna}{%
  \rotatebox[origin=c]{\despachoSelloAngulo}{%
    \includegraphics[height=\despachoSelloTam]{\despachoSelloImagen}%
  }%
}

%% --- Bloque de firma --------------------------------------------------
%%  Firma \meta{firmante}/\meta{cargofirmante} si existen (personal del
%%  despacho); si no, el titular. El sello va al costado derecho via
%%  \rlap (caja de ancho cero: no altera el centrado de la firma).
\newcommand{\bloquefirma}[1][2.2cm]{%
  \vspace{#1}%
  \noindent
  \ifdespachosello
    % Overlay de dimension CERO (\rlap + raisebox[0pt][0pt]) que
    % CUELGA bajo la linea de firma: el sello queda debajo del cargo,
    % alineado al margen DERECHO, sin consumir espacio del flujo (una
    % caja real empujaria la firma de pagina en documentos justos).
    % Caida = alto del arte + separacion + media altura del bloque.
    \rlap{\makebox[\textwidth][r]{%
      \raisebox{\dimexpr-\despachoSelloTam-\despachoSelloAlt-0.50cm\relax}[0pt][0pt]{\selloconsigna}%
      \hspace*{\despachoSelloPos}}}%
  \fi
  \makebox[\textwidth][c]{%
    % minipage = caja inquebrantable: la firma nunca se parte entre paginas
    \begin{minipage}[c]{0.6\textwidth}\centering
      {\seriesemi\MakeUppercase{\ifmetadef{firmante}{\meta{firmante}}{\meta{autor}}}}\\[1pt]
      \ifmetadef{cargofirmante}{\meta{cargofirmante}}{\meta{cargo}}%
    \end{minipage}%
  }%
  \ifdespachosello
    % Puntal invisible: reserva bajo la linea la profundidad que ocupa
    % el sello colgante MENOS 0.90cm de tolerancia (el sello puede
    % asomarse al margen inferior sin tocar el pie, que empieza a
    % footskip=1.10cm). Si aun asi no cabe, la firma y su sello pasan
    % completos a la pagina siguiente en vez de pisar el pie.
    \rule[\dimexpr-\despachoSelloTam-\despachoSelloAlt-0.50cm+0.90cm\relax]{0pt}{0pt}%
  \fi
  \par
}

%% --- Sello anclado al pie de la pagina de la firma --------------------
%%  Posicion FIJA del sello, abajo a la derecha de la pagina donde cae la
%%  firma, sin ocupar espacio del texto (Edison 2026-09-07). Dos casos
%%  (Edison 2026-10-01: «el sello debe estar justo ENCIMA de la nota de
%%  pagina, no debajo»):
%%
%%    · pagina SIN notas al pie: el sello vive en la BANDA DEL PIE, apoyado
%%      sobre el filete del pie institucional (como hasta ahora);
%%    · pagina CON notas al pie: el sello va JUSTO ENCIMA de las notas,
%%      alineado a la derecha, y las notas bajan a la banda del pie el alto
%%      que el sello deja libre. Es un INTERCAMBIO: el texto no se mueve, la
%%      paginacion no cambia y el sello nunca queda debajo de una nota.
%%
%%     \sinsello              % en el preambulo: no cuelga de la firma
%%     ...
%%     \bloquefirma[2.0cm]
%%     \selloalfondo          % sello al pie de la pagina de la firma
%%
%%  Como funciona: \selloalfondo deja en el flujo una INSERCION vacia de la
%%  clase \despachoInsSello (alto cero, \count 0: no reserva nada). Una
%%  insercion viaja con SU pagina, asi que la rutina de salida sabe con
%%  exactitud que pagina lleva la firma (antes se usaba «la proxima pagina
%%  que salga» de eso-pic, que con el constructor de paginas atrasado podia
%%  equivocarse de hoja). Antes de armar la columna (\@makecol) se mira si
%%  esa pagina trae notas (\footins):
%%    - sin notas: overlay de eso-pic sobre ESTA pagina, en la banda del pie
%%      (base = filete del pie + \despachoSelloSobreFilete);
%%    - con notas: el sello se antepone a la caja de notas, seguido de la
%%      raya de notas, y la raya original se apaga solo para esa pagina.
%%
%%    \setsellosobrefilete{0.12cm}  aire entre el filete y el sello
%%    \setsellosobrenotas{0.25cm}   aire entre el sello y la raya de notas
%%    \setselloajuste{-0.3cm}       ajuste EXCEPCIONAL: negativo = sube,
%%                                  positivo = baja (panel Formato del SGDP)
\newlength{\despachoSelloSobreFilete}
\setlength{\despachoSelloSobreFilete}{0.12cm}
\newcommand{\setsellosobrefilete}[1]{\setlength{\despachoSelloSobreFilete}{#1}}
\newlength{\despachoSelloSobreNotas}
\setlength{\despachoSelloSobreNotas}{0.25cm}  % 2026-10-12: 0.10 dejaba el sello rozando la 1.ª linea de la nota
\newcommand{\setsellosobrenotas}[1]{\setlength{\despachoSelloSobreNotas}{#1}}
\newlength{\despachoSelloAjuste}
\setlength{\despachoSelloAjuste}{0pt}
\newcommand{\setselloajuste}[1]{\setlength{\despachoSelloAjuste}{#1}}
\RequirePackage{etoolbox}   % \pretocmd
\newinsert\despachoInsSello
\count\despachoInsSello=0
\dimen\despachoInsSello=\maxdimen
\skip\despachoInsSello=\z@
\newcommand{\selloalfondo}{%
  % respeta el interruptor de documento (\setmeta{sello}{no})
  \ifdespachosellodoc
    \insert\despachoInsSello{\hrule\@height\z@\@width\z@}%
  \fi
}
%%  Overlay en la banda del pie de la pagina que se esta emitiendo.
\newcommand{\despachoSelloEnBanda}{%
  \AddToShipoutPictureFG*{%
    \AtPageLowerLeft{%
      % x: borde derecho del area de texto (menos la separacion opcional)
      \hspace*{\dimexpr 1in+\oddsidemargin+\textwidth-\despachoSelloPos\relax}%
      % y: base del sello = filete del pie + aire (- ajuste = sube)
      \raisebox{\dimexpr\paperheight-1in-\topmargin-\headheight-\headsep-\textheight
                        -\footskip+\despachoSelloSobreFilete-\despachoSelloAjuste\relax}%
               [\z@][\z@]{\makebox[\z@][r]{\selloconsigna}}%
    }%
  }%
}
%%  Raya de notas del documento, guardada al empezar (una clase pudo
%%  cambiarla) para dibujarla DEBAJO del sello en la pagina que lo lleva.
\AtBeginDocument{\let\despachoRayaNotas\footnoterule}
%%  Aire minimo entre la ultima linea del texto (la firma) y la raya de notas en la hoja del sello
%%  (Edison 2026-10-14, oficio 144 compactado a una hoja: «Diputado de la Republica» tocaba la raya;
%%  «siempre debe estar ligeramente separado»). Si falta, las notas bajan lo necesario dentro de la
%%  banda del pie, como mucho \despachoNotasExtraMax, para no pisar el filete.
%%     \setnotasairemin{0.30cm}
\newdimen\despachoSelloSobreRaya
\newdimen\despachoTmpAire
\newlength{\despachoNotasAireMin}
\setlength{\despachoNotasAireMin}{0.30cm}
\newcommand{\setnotasairemin}[1]{\setlength{\despachoNotasAireMin}{#1}}
\newlength{\despachoNotasExtraMax}
\setlength{\despachoNotasExtraMax}{0.35cm}
\newcommand{\despachoSelloAntesDeColumna}{%
  \ifvoid\despachoInsSello\else
    \global\setbox\despachoInsSello\box\voidb@x
    % Dos lugares fijos, nunca encima de una nota (Edison 2026-10-12, con foto: «el sello esta
    % debajo de la nota de pie; deberia estar mas arriba, sin superponerse; con notas, ese sello
    % sube arriba»). Restituye la regla del 2026-10-01:
    %   - hoja SIN notas: el sello en la banda del pie (junto a Adj. y c.c. si los hay, que es el
    %     unico «aprovechar el espacio» pedido, y sin superponerse);
    %   - hoja CON notas: el sello JUSTO ENCIMA de las notas, a la derecha, y la raya de notas
    %     debajo de el.
    % La variante «sello en la banda aunque haya notas» (2026-10-10, segunda vuelta) y la de «sello
    % al lado de las notas cortas» (manana del 2026-10-10) lo ponian SOBRE las notas, que bajan a esa
    % banda (\despachoNotasBajan): se retiran.
    \ifvoid\footins
      \despachoSelloEnBanda
    \else
      % Todo es local a la rutina de salida: la pagina siguiente vuelve a la raya normal.
      \setbox\tw@\hbox{\selloconsigna}%
      % Lo que hay encima de la raya (ajuste + sello + aire): lo usa el gancho de las notas para
      % dejar siempre un aire minimo entre la firma y la raya (2026-10-14).
      \global\despachoSelloSobreRaya\dimexpr\despachoSelloAjuste+\ht\tw@+\dp\tw@+\despachoSelloSobreNotas\relax
      \setbox\footins\vbox{%
        \kern\despachoSelloAjuste
        \hbox to\columnwidth{\hfill\box\tw@\hspace*{\despachoSelloPos}}%
        \vskip\despachoSelloSobreNotas
        \despachoRayaNotas
        \unvbox\footins}%
      \let\footnoterule\relax
    \fi
    \global\despachoNotaAncho\z@
  \fi
}
%% --- Notas al pie mas abajo (2026-10-06) -------------------------------
%%  Bajo las notas quedaba ~0.9 cm en blanco hasta el filete del pie: la
%%  banda del pie (footskip=1.50cm) se ensancho para el sello, que en una
%%  pagina con notas va ENCIMA de ellas (Edison 2026-10-06: «hay mucho
%%  espacio debajo de la nota de pie; bajarla un poco mas de la mitad»).
%%  Al armar la columna, las notas (raya incluida) bajan
%%  \despachoNotasBajan dentro de esa banda. Solo se mueven al componer la
%%  pagina: el corte de pagina ya esta decidido, asi que la paginacion y el
%%  texto no cambian. Local a la rutina de salida.
%%     \setnotasbajan{1.32cm}   (bajo el borde del texto; deja ~0,4 cm hasta el filete)
\newlength{\despachoNotasBajan}
\setlength{\despachoNotasBajan}{1.32cm}
\newcommand{\setnotasbajan}[1]{\setlength{\despachoNotasBajan}{#1}}
%%  Y SIEMPRE al fondo de la hoja (2026-10-06, OF-0123): con \raggedbottom
%%  las notas subian pegadas al texto cuando la hoja no se llenaba (el
%%  cierre pasaba a la siguiente) y el blanco quedaba DEBAJO de ellas. Al
%%  armar la columna se rellena con lo que le falta al texto para la altura
%%  de la hoja, como la opcion «bottom» de footmisc.
%%  La altura se calcula sobre la hoja entera (\@colht): texto natural +
%%  separacion + notas (con el sello, si va encima de ellas) = alto de la
%%  hoja; lo que falta va entre el texto y la raya. Se engancha ANTES que
%%  el sello (\pretocmd antepone: el ultimo enganchado corre primero), asi
%%  que ya cuenta el sello que \despachoSelloAntesDeColumna puso en \footins.
%%  Y el TEXTO aprovecha ese espacio (2026-10-06, OF-0123/0124: el cierre
%%  pasaba a una tercera hoja con blanco de sobra en la segunda): la
%%  separacion texto-notas que mide el corte de pagina se descuenta en
%%  \despachoNotasBajan, porque las notas viven en la banda del pie. Al armar
%%  la columna se repone: entre texto y notas queda al menos la separacion
%%  normal. Despues de todo el preambulo, para respetar el \setnotasbajan de
%%  cada clase.
\AfterEndPreamble{\addtolength{\skip\footins}{-\despachoNotasBajan}}
\pretocmd{\@makecol}{%
  \ifvoid\footins\else
    \setbox\z@\vbox{\unvcopy\@cclv}%
    \dimen@\dimexpr\@colht+\despachoNotasBajan-\ht\z@-\dp\z@-\skip\footins-\ht\footins-\dp\footins\relax
    % Negativo solo en la hoja del sello (va encima de las notas): el texto
    % reserva bajo la firma el hueco del sello colgante (puntal de \bloquefirma),
    % que aqui queda vacio y las notas suben sobre el. En otra hoja, nunca.
    \ifx\footnoterule\relax
      % (sin tope fijo: si las notas no suben, la medida fija se pierde y pisan el pie). Pero
      % nunca pegadas a la firma (2026-10-14): entre el texto y la raya queda al menos
      % \despachoNotasAireMin; para lograrlo las notas bajan como mucho \despachoNotasExtraMax.
      \despachoTmpAire\dimexpr\despachoNotasAireMin-\despachoSelloSobreRaya-\skip\footins\relax
      \ifdim\despachoTmpAire>\dimexpr\dimen@+\despachoNotasExtraMax\relax
        \despachoTmpAire\dimexpr\dimen@+\despachoNotasExtraMax\relax
      \fi
      \ifdim\dimen@<\despachoTmpAire\dimen@\despachoTmpAire\fi
    \else
      \ifdim\dimen@<\z@\dimen@\z@\fi
    \fi
    \advance\skip\footins\dimen@
  \fi}{}{%
  \ClassWarning{despacho}{No se pudo enganchar las notas a \string\@makecol}}

\pretocmd{\@makecol}{\despachoSelloAntesDeColumna}{}{%
  \ClassWarning{despacho}{No se pudo enganchar el sello a \string\@makecol}}

%% --- Notas al pie SIN sangria (2026-10-06) ----------------------------
%%  Edison: «las notas al pie no deben tener sangria». La llamada (¹) y el
%%  texto empiezan al margen, en la primera linea y en las demas. Una clase
%%  puede redefinir \@makefntext despues (proyecto-de-ley tiene la suya).
%%  (El ancho natural de cada nota se sigue midiendo en \despachoNotaAncho, pero desde el
%%  2026-10-12 el sello ya no se pone al lado de las notas: con notas va siempre encima.)
\newdimen\despachoNotaAncho
\global\despachoNotaAncho\z@
\renewcommand\@makefntext[1]{%
  \begingroup\despachotamnotas
    \sbox\z@{\despachotamnotas #1}%
    \ifdim\wd\z@>\despachoNotaAncho\global\despachoNotaAncho\wd\z@\fi
  \endgroup
  \parindent\z@\noindent{\despachotamnotas\@makefnmark\,#1}}

%% --- Espacio ANTES DEL CUERPO (2026-10-01) ----------------------------
%%  El aire entre la cabecera —asunto o sumilla, referencia— y el primer
%%  parrafo (Edison 2026-10-01: «debe poder controlarse el espacio entre la
%%  Sumilla y el primer parrafo»). Factor de \baselineskip, como los demas
%%  espacios de cabecera, para que siga al interlineado. Vacio = el de cada
%%  estructura (oficio: el mismo de las partes, 0.20; memorando e informe:
%%  0.8 bajo el filete). Lo consumen las estructuras:
%%     \espacioantescuerpo{0.8}          (memorando, informe)
%%     \despachoAjusteAntesSaludo        (oficio: corrige su \espaciopartes)
\newcommand{\despachoEspacioCuerpo}{}
\newcommand{\setespacioantescuerpo}[1]{\renewcommand{\despachoEspacioCuerpo}{#1}}
\newcommand{\espacioantescuerpo}[1]{%
  \vspace{\ifx\despachoEspacioCuerpo\@empty #1\else\despachoEspacioCuerpo\fi\baselineskip}%
}

%% --- Pie: adjuntos y copias ------------------------------------------
%% --- Pie con medidas propias (Edison 2026-10-10: «controlar por separado la posicion, el
%%  tamano de letra y el espaciado de la firma, las notas y la lista de adjuntos») -------
%%    \settampie{\fontsize{9}{10.8}\selectfont}   letra de Adj. (e iniciales)
%%    \setespaciopie{0.4}                          aire sobre Adj., en lineas (defecto 0,8)
%%    \settamnotas{\fontsize{8}{9.6}\selectfont}  letra de las notas al pie
\newcommand{\despachotampie}{\small}
\newcommand{\settampie}[1]{\renewcommand{\despachotampie}{#1}}
\newcommand{\despachoEspacioPie}{0.8}
\newcommand{\setespaciopie}[1]{\renewcommand{\despachoEspacioPie}{#1}}
\newcommand{\despachotamnotas}{}
\newcommand{\settamnotas}[1]{\renewcommand{\despachotamnotas}{#1}}

\newcommand{\adjunto}[1]{%
  \vspace{\despachoEspacioPie\baselineskip}%
  \noindent{\despachotampie \textbf{Adj.:} #1\par}%
}
\newcommand{\concopia}[1]{%
  \vspace{0.3\baselineskip}%
  \noindent{\ifx\despachotamcopia\@empty\small\else\despachotamcopia\fi
    \despachointercopia \textbf{c.c.:} #1\par}%
}

%% --- Pie en LISTA: adjuntos y copias, uno por linea (2026-10-01) ------
%%  «Los adjuntos se muestran de manera horizontal y se ve feo»; «la c.c.
%%  ocupa demasiado». Lista formal y compacta: la etiqueta a la izquierda y
%%  cada elemento en su linea, alineados; uno largo sangra bajo su propio
%%  texto. \adjunto y \concopia (en una linea) siguen para los .tex a mano.
%%
%%     \listapie{\small}{Adj.:}{\itempie[1.]{Informe tecnico}\itempie[2.]{Acta}}
%%     \listapie{\small}{c.c.:}{\itempie{Nombre Apellido, cargo}}
\newlength{\despachoEtiquetaPie}
\newlength{\despachoMarcaPie}
\newcommand{\itempie}[2][]{%
  \if\relax\detokenize{#1}\relax
    {\hangindent=1em\hangafter=1 \noindent #2\par}%
  \else
    \settowidth{\despachoMarcaPie}{#1\ }%
    {\hangindent=\despachoMarcaPie\hangafter=1 \noindent #1\ #2\par}%
  \fi
}
\newcommand{\listapie}[3]{%
  \par\noindent{#1%
    % Columna comun para «Adj.:» y «c.c.:» (la etiqueta mas ancha), asi
    % las dos listas quedan alineadas una bajo otra.
    \settowidth{\despachoEtiquetaPie}{\textbf{Adj.:}\ }%
    \settowidth{\despachoMarcaPie}{\textbf{#2}\ }%
    \ifdim\despachoMarcaPie>\despachoEtiquetaPie
      \setlength{\despachoEtiquetaPie}{\despachoMarcaPie}\fi
    \makebox[\despachoEtiquetaPie][l]{\textbf{#2}}%
    \parbox[t]{\dimexpr\linewidth-\despachoEtiquetaPie\relax}{%
      \setlength{\parskip}{0pt}\setlength{\parindent}{0pt}#3}%
    \par}%
}
\newcommand{\adjuntos}[1]{%
  \vspace{\despachoEspacioPie\baselineskip}%
  \listapie{\despachotampie}{Adj.:}{#1}%
}
\newcommand{\concopias}[1]{%
  \vspace{0.3\baselineskip}%
  \listapie{\ifx\despachotamcopia\@empty\small\else\despachotamcopia\fi\despachointercopia}{c.c.:}{#1}%
}

%% Iniciales de responsabilidad: ultima linea del pie, en versalita pequena
%% y al margen izquierdo («ABC/xyz»: quien firma / quien redacta o digita).
%% Es una PARTE del oficio administrativo —identifica al responsable de la
%% redaccion— que el sistema no componia (2026-09-23).
\newcommand{\iniciales}[1]{%
  \vspace{0.3\baselineskip}%
  \noindent{\footnotesize\despachotampie #1\par}%
}

%% --- Tamaño de letra del bloque destinatario (ajustable) -------------
%%  \settamdestinatario{\small}  (o \footnotesize) en el preambulo reduce
%%  el bloque de destinatario. Por defecto, tamaño normal del cuerpo. Va
%%  aqui (compartido) para que el comando exista en TODOS los tipos; solo
%%  lo consumen las estructuras que tienen destinatario (oficio, solicitud).
\newcommand{\despachotamdest}{}
\newcommand{\settamdestinatario}[1]{\renewcommand{\despachotamdest}{#1}}

%% --- Ajustes FINOS por parte (Edison 2026-09-07) ---------------------
%%  Interlineado del bloque destinatario y tamano/interlineado de las
%%  copias (arriba, dentro de \destinatario, o al pie con \concopia),
%%  independientes del cuerpo. Los emite el generador del SGDP (campos
%%  fmt_dest_interlineado, fmt_copia_tamano, fmt_copia_interlineado).
%%    \setinterlineadodest{1.0}     \settamcopia{\small}
%%    \setinterlineadocopia{0.9}
\newcommand{\despachointerdest}{}
\newcommand{\setinterlineadodest}[1]{\renewcommand{\despachointerdest}{\setstretch{#1}}}
\newcommand{\despachotamcopia}{}
\newcommand{\settamcopia}[1]{\renewcommand{\despachotamcopia}{#1}}
\newcommand{\despachointercopia}{}
\newcommand{\setinterlineadocopia}[1]{\renewcommand{\despachointercopia}{\setstretch{#1}}}

% --- Fotografía insertada desde el editor del SGDP (2026-10-05) ---
%  \fotografia{archivo}{ancho}{título}{fuente}: imagen centrada (ancho como
%  fracción del texto, nunca más de media página de alto), rótulo numerado
%  «FOTOGRAFÍA — N° n» y título debajo; la fuente va a nota al pie. Sin
%  flotante: queda en el lugar del texto en que se insertó.
\newcounter{fotografia}
\newcommand{\fotografia}[4]{%
  \par\medskip
  \begingroup\centering
    \includegraphics[width=#2\linewidth,height=0.5\textheight,keepaspectratio]{#1}\par
    \smallskip\refstepcounter{fotografia}%
    {\color{franjacongreso}\bfseries FOTOGRAFÍA — N° \thefotografia}\par
    {\color{franjacongreso}\itshape #3}%
    \if\relax\detokenize{#4}\relax\else\footnote{#4}\fi\par
  \endgroup
  \medskip}
